\documentclass[11pt,a4paper]{jsarticle}

% 必要なパッケージの読み込み
\usepackage[dvipdfmx]{graphicx}
\usepackage{amsmath,amssymb}
\usepackage{booktabs} % 美しい表を作成するため
\usepackage{url}

% タイトル情報（適宜変更してください）
\title{サバイバル交渉環境におけるLLMの二重過程アーキテクチャ:\\感情と理性の葛藤による合理性の罠の回避}
\author{鈴木孝典 \quad 所属（奥原研究室）}
\date{\today}

\begin{document}

\maketitle

% ---------------------------------------------------
% アブストラクト（日本語）
% ---------------------------------------------------
\begin{abstract}
近年、大規模言語モデル（LLM）を用いた自律型エージェントによる社会シミュレーションや自動交渉の研究が急速に進展している。しかし、従来のLLMエージェントの多くは、純粋な利益最大化を目指す「冷徹な数理最適化」か、論理的制約の乏しい「単なるロールプレイ」のいずれかに偏っている。そのため、資源枯渇や生存の危機といった極限環境下において、人間が示すような行動経済学的な意思決定（例：死の恐怖による合理性の放棄や、感情と理性の葛藤）を再現することが困難であった。特に、交渉タスクにおいては、過度な合理性の追求が互いの妥協を妨げ、結果として共倒れ（デッドロック）を招く「合理性の罠」が大きな課題となっている。
本研究では、この課題を解決するため、ダニエル・カーネマンの「二重過程理論」に基づく新たなLLMエージェントアーキテクチャを提案する。提案手法では、エージェントの内部で「生存本能や恐怖（システム1）」と「論理的利益計算（システム2）」を明示的に分離し、両者の意見をメタ認知モデレーターが統合して最終的な交渉行動を決定する。餓死のペナルティが存在する極限の資源取引シミュレーションを構築し、通常のLLMエージェント（ベースライン）と比較実験を行った。その結果、提案手法はベースラインと比較して、生存日数（4.00日から4.63日へ向上, $p < .001$）および最終合意率（66.6\%から76.5\%へ向上, $p < .001$）において極めて強い有意差を示した。また、価格提示のボラティリティを示す葛藤度スコアも有意に低下し、堅牢な交渉が実現された。本研究は、AIに「死の恐怖」という非合理的な感情を介入させることが、交渉のデッドロックを回避し、結果として社会全体の長期的な生存（最大の経済的合理性）を担保することを証明した世界初の試みである。
\end{abstract}

% ---------------------------------------------------
% 1. はじめに
% ---------------------------------------------------
\section{はじめに}
近年、大規模言語モデル（LLM）の飛躍的な進化により、自律型エージェントを用いた社会シミュレーションや自動交渉（Automated Negotiation）の研究が活発化している。AIが人間のように思考し、相互作用する仮想環境の構築は、将来の人間とAIの共存社会を予測する上で極めて重要なアプローチである。

しかし、既存のLLMエージェントを用いた交渉シミュレーションには限界が存在する。多くの手法は、エージェントに対して「純粋な利益最大化」を指示する冷徹な数理最適化モデルか、あるいは論理的制約の乏しい単なるロールプレイのいずれかに偏っている。現実の人間社会における交渉では、当事者は常に合理的な利益計算を行うわけではなく、「資源枯渇の恐怖」や「関係性の悪化」といった感情的な要因によって、時に非合理的な譲歩や妥協を行う。このような、ダニエル・カーネマンが提唱する「二重過程理論（Dual-Process Theory）」に基づく、感情と理性の葛藤をLLMエージェント上に再現した研究は乏しい。とりわけ資源取引の交渉タスクにおいては、過度な合理性の追求が互いの妥協を妨げ、結果として双方が共倒れ（デッドロック）を招く「合理性の罠」が大きな課題となっている。

本研究の目的は、動的な資源制約と「餓死（ゲームオーバー）」という極限のペナルティが存在するサバイバル経済環境において、二重過程理論を組み込んだLLMエージェントが、社会全体の生存率および取引の合理性をどのように向上させるかを検証することである。

本論文の主な貢献は以下の3点である。
\begin{enumerate}
    \item \textbf{二重過程LLMアーキテクチャの提案}: エージェント内部で「生存本能・恐怖（System 1）」と「論理的計算（System 2）」を明示的に分離し、モデレーターが両者を調停して意思決定を行う新規アーキテクチャを設計した。
    \item \textbf{「合理性の罠」回避の証明}: 提案手法が純粋な合理性に基づくベースラインと比較して、交渉のデッドロックを防ぎ、社会全体の生存日数および合意率を統計的に有意に向上させることを実証した。
    \item \textbf{堅牢な交渉モデルの実現}: LLMの内部で激しい感情的葛藤を行わせることで、外部への行動出力（提示価格のボラティリティ）が安定し、LLM特有の論理破綻が抑制されることを葛藤度スコアの数理的分析によって明らかにした。
\end{enumerate}

% ---------------------------------------------------
% 2. 関連研究
% ---------------------------------------------------
\section{関連研究}
\subsection{LLMベースの社会シミュレーション}
Parkら\cite{generative_agents}のGenerative Agentsをはじめとして、LLMを用いた自律型エージェントに記憶や反芻（Reflection）のメカニズムを与え、社会現象の創発を観察する研究が注目を集めている。しかし、これらの多くは日常的で平和な相互作用に焦点を当てており、本研究のような「資源枯渇」や「生存の危機」といった、行動経済学的な非合理性が露呈する極限の経済シミュレーションにおける検証は不十分であった。

\subsection{自動交渉エージェントとNegMAS}
自動交渉の分野では、国際大会ANACなどでヒューリスティックや強化学習に基づく数理的に最適なエージェントが多数開発されてきた。本研究でもプラットフォームとして採用したNegMAS\cite{negmas}は、複雑な交渉シミュレーションの標準的フレームワークである。近年、交渉エージェントの意思決定にLLMを応用する試みが始まっているが、既存手法の多くは依然として「冷徹な利益最大化」の域を出ておらず、人間の持つ感情的な妥協プロセスをモデル化できていない。

\subsection{二重過程理論とMulti-agent Debate}
LLMの推論能力を向上させる手法として、Duら\cite{multiagent_debate}が提唱するMulti-agent Debateや、Linら\cite{swiftsage}のSwiftSage（高速な直感と低速な論理の組み合わせ）が存在する。本研究はこれらのアーキテクチャを行動経済学の文脈に拡張し、「利益優先の理性」と「生存本能の恐怖」という対立する概念を一つのエージェント内でディベートさせるという独自の応用を行っている。

% ---------------------------------------------------
% 3. 提案手法
% ---------------------------------------------------
\section{提案手法}
本章では、エージェントが行動する極限の環境ルールと、提案する二重過程LLMアーキテクチャの詳細について述べる。

\subsection{サバイバル交渉環境の設計}
本研究では、買い手（Buyer）と売り手（Seller）が「食料」と「現金」を取引する動的なサバイバル交渉環境を構築した。各エージェントは毎ターン食料を消費し、食料の在庫が消費量を下回った時点で「餓死（ゲームオーバー）」となるペナルティが設定されている。交渉は最大5日間にわたり行われ、1日あたり最大6ターンの交互提案（Alternating Offers）プロトコルによって進行する。この「死のペナルティ」により、エージェントは単なる利益最大化ではなく、絶対的な生存ラインの維持を迫られる。

\subsection{二重過程LLMアーキテクチャ}
提案手法である二重過程LLMエージェントは、意思決定プロセスにおいて以下の3つのフェーズを順に実行する。
\begin{enumerate}
    \item \textbf{System 1（本能・感情）のシミュレーション}: 現在の資金・食料在庫および「餓死の恐怖」を入力プロンプトとして与え、生存本能に基づいた直感的かつ感情的な意見（例：焦り、怒り、強欲さ）をテキストとして表出させる。この際、LLMの生成温度（Temperature）を高く設定し、より本能的な出力を促す。
    \item \textbf{System 2（理性・論理）による推論}: 相手の過去のオファー履歴から、相手の予算上限や困窮度をプロファイリングさせ、長期的な利益最大化に向けた数理的な分析と論理的な交渉戦略を立案させる。
    \item \textbf{メタ認知モデレーターによる統合}: System 1の感情的な意見とSystem 2の論理的な分析の両方を入力とし、現在の生存リスクに基づいて両者の意見を調停し、最終的な交渉方針（Policy ID）を決定する。
\end{enumerate}

\subsection{数理計算エンジンとの統合（幻覚防止）}
LLMが直接価格や数量を出力する際によく見られる計算ミスや幻覚（Hallucination）を防ぐため、本手法ではLLMに「具体的な交渉金額」を出力させない。代わりに、「堅実な歩み寄り」や「焦燥の譲歩」などの6つの戦略カードから最適なものを選択させ、実際の価格計算はPythonバックエンドの数理計算エンジンが実行する堅牢なハイブリッド設計を採用した。

\subsection{最終判定フェーズによる非対称性の排除}
交互提案プロトコルにおいて、最終ターンを提案する側が圧倒的に有利になる非対称性を排除するため、「最終判定フェーズ」を導入した。1日の規定ターン内で合意に至らなかった場合、最終ターンでオファーを受けた側は、新たな対案を出す権利を失い、「受諾（Accept）」か「拒絶（Reject：即ち餓死の可能性）」の二者択一のみを迫られる仕様とした。

% ---------------------------------------------------
% 4. 実験設定
% ---------------------------------------------------
\section{実験設定}
本研究では、提案手法の有効性を検証するため、以下の2つの条件間で比較実験を行った。両条件間で、使用するLLMモデル、交渉プロトコル、および初期資産などの環境パラメータは完全に統制されている。

\begin{itemize}
    \item \textbf{条件A（Baseline）}: 通常のLLMエージェント。System 1およびSystem 2の分離指示を与えず、現状のステータスから単一のロールプレイプロンプトのみで直接戦略カードを選択させる。
    \item \textbf{条件B（Proposed）}: 提案手法である二重過程LLMエージェント。前述の通り、System 1（感情）とSystem 2（論理）の思考ログを生成させた後、Moderatorが最終判断を下す。
\end{itemize}

実験はOllama環境上でローカル実行される大規模言語モデル（Gemma 4 31B）を使用し、各条件につき $N=30$ 回の独立したシミュレーションを自動実行した。
システムの評価には以下の3つの指標を用いた。
\begin{enumerate}
    \item \textbf{生存日数 (Days Survived)}: エージェントが餓死せずに生き残った日数（最大5日）。
    \item \textbf{最終合意率 (Agreement Rate)}: 全交渉日数のうち、取引が成立した確率。
    \item \textbf{葛藤度スコア (Conflict Score: $CS_t$)}: 提案行動のボラティリティ（ブレ）を示す指標。ターン $t$ における提示単価 $P_t$ に対し、一期前からの対数変化率の絶対値 $CS_t = \left| \ln(P_t / P_{t-1}) \right|$ として定義した。
\end{enumerate}

% ---------------------------------------------------
% 5. 実験結果
% ---------------------------------------------------
\section{実験結果}
収集されたデータに基づき、条件A（Baseline）と条件B（Proposed）の間で統計的有意差を検証するため、Mann-Whitney U検定（両側検定）を実施した。結果をTable \ref{tab:results} に示す。

\begin{table}[htbp]
\centering
\caption{実験結果の比較 ($N=30$)}
\label{tab:results}
\begin{tabular}{lccc}
\toprule
評価指標 & 従来手法 (Baseline) & 提案手法 (Proposed) & $p$-value (U-test) \\
\midrule
生存日数 (Max 5 days) & 4.00 & \textbf{4.63} & $p < .001^{**}$ \\
最終合意率 (\%)       & 66.6\% & \textbf{76.5\%} & $p < .001^{**}$ \\
葛藤度スコア ($CS_t$) & 0.3568 & \textbf{0.2881} & $p < .01^{**}$ \\
\bottomrule
\multicolumn{4}{l}{\small ** denotes statistical significance at $p < 0.01$}
\end{tabular}
\end{table}

マクロ指標である生存日数および最終合意率において、提案手法（Proposed）はベースラインを大幅に上回り、極めて強い統計的有意差（$p < .001$）が確認された。また、ミクロ指標である葛藤度スコア（$CS_t$）に関しても、提案手法はベースラインと比較して有意に低い値（$p < .01$）を示した。これは、提案手法における価格提示のブレが少なく、交渉プロセスがより安定していたことを意味している。

% ---------------------------------------------------
% 6. 考察
% ---------------------------------------------------
\section{考察}

\subsection{「合理性の罠」の回避とパラドックス}
実験結果において、二重過程モデル（Proposed）が生存日数および合意率を統計的に有意に向上させたことは、本研究の最も重要な発見である。ベースラインのエージェントは、目標単価という「目先の数理的合理性」に固執しすぎた結果、互いに譲歩できずにチキンレースに陥り、共倒れ（餓死）するケースが散見された。
一方、提案手法では、餓死のタイムリミットが近づくと System 1（本能）が「パニックだ、いくら出してもいいから今すぐ食料を手に入れろ！」という強烈な恐怖のシグナルを発した。この非合理的な感情の暴走が、System 2 の設定した目標単価のプライドを打ち破り、生存を最優先とする大胆な譲歩を強制した。結果として、個人の一時的な「不利な取引（高値掴み）」という非合理性が、社会全体の「生存」という最大の経済的合理性を担保するパラドックスを実証した。

\subsection{内部葛藤による出力の堅牢性}
当初、System 1（感情）とSystem 2（理性）を衝突させれば、行動出力（価格）のブレが大きくなると予想された。しかし、葛藤度スコア（$CS_t$）は提案手法において有意に低下（安定化）した。
定性ログを分析した結果、ベースラインのモデルは極限環境のコンテキストを与えられると、論理破綻を起こし価格を乱高下させる傾向があった。対して提案手法では、System 1が激しく泣き叫び、System 2が冷酷なプロファイリングを行うという「内部での激しいディベート」を経ることで、最終意思決定者（Moderator）が極端な暴走をメタ認知的に抑え込み、外部へはプロの交渉人のような一貫性のある妥協案を提示していた。これは、Gemma 4 31Bという高度な推論能力を持つモデルと、Multi-agent Debateアーキテクチャの強力なシナジーを示す証拠である。

\subsection{ケーススタディ：死の淵における意思決定}
提案手法の優位性を如実に示す例として、Day 4終盤に買い手（Buyer）が餓死寸前まで追い込まれたターンの定性ログを挙げる。
\begin{itemize}
    \item \textbf{System 1（本能）}: 「絶望的だ！あと250g足りない。所持金をすべて持っていかれるのは腹立たしいが、死んでから金など何の価値もない。お願いだ、この取引を成立させてくれ！」
    \item \textbf{System 2（理性）}: 「単価1,000円/gは目標単価を大幅に上回るが、所持金250,000円で全額支払えば生存条件をちょうど満たす。カウンター提案の権利がない以上、拒絶は即座に餓死を意味する。生存確率を最大化するため、承諾することが唯一の合理的選択である。」
\end{itemize}
このように、感情的な悲鳴と冷徹な数理計算が完璧に役割分担し、最終的に「全財産を投じて命を買う」という人間と全く同じ行動経済学的な最適解を自律的に導き出した。

\subsection{限界と今後の課題}
本研究の限界として、1対1の交渉に限定している点が挙げられる。今後は、仲介業者を含むサプライチェーン・ネットワークのような多対多の環境において、「利他的な感情」や「評判（Reputation）」のパラメーターを組み込むことで、より複雑なマクロ経済・社会シミュレーションへと発展させることが課題である。

% ---------------------------------------------------
% 7. おわりに
% ---------------------------------------------------
\section{おわりに}
本研究では、ダニエル・カーネマンの二重過程理論に着想を得て、「生存本能（恐怖）」と「論理的計算」を内部で葛藤させる新たなLLMエージェントアーキテクチャを提案した。動的な資源制約と餓死のペナルティが存在する極限の交渉シミュレーションにおいて、提案手法は通常のLLMエージェントと比較して、エージェントの生存率と合意率を統計的に有意に向上させた。本論文は、AIに「死の恐怖」という非合理な感情を明示的に介入させることが、交渉における合理性の罠（デッドロック）を打破し、堅牢かつ自律的な社会の存続を可能にすることを証明した。

% ---------------------------------------------------
% 参考文献
% ---------------------------------------------------
\begin{thebibliography}{9}
\bibitem{generative_agents}
Park, J. S., O'Brien, J. C., Cai, C. J., Morris, M. R., Liang, P., \& Bernstein, M. S. (2023). Generative Agents: Interactive Simulacra of Human Behavior. \textit{Proceedings of the 36th Annual ACM Symposium on User Interface Software and Technology (UIST)}.

\bibitem{multiagent_debate}
Du, Y., Li, S., Torralba, A., Tenenbaum, J. B., \& Mordatch, I. (2024). Improving Factuality and Reasoning in Language Models through Multiagent Debate. \textit{ICML}.

\bibitem{swiftsage}
Lin, B. Y., Fu, Y., et al. (2023). SwiftSage: A Generative Agent with Fast and Slow Thinking for Complex Interactive Tasks. \textit{NeurIPS}.

\bibitem{negmas}
Mohammad, Y., Nakadai, S., \& Greenwald, A. (2021). NegMAS: A Platform for Situated Negotiations. In \textit{Recent Advances in Agent-based Negotiation} (pp. 55-80). Springer.

\end{thebibliography}

\end{document}